% BHU PYQ -- Transcribed & Corrected
% Paper: CSCSE31: Python Programming (Sessional Exam)
% Exam: B.Sc. Semester III Sessional Examination 2025-26 (NEP)
% Category: Skill Enhancement Course (SEC)
% Department: Computer Science

\documentclass[11pt,a4paper]{article}
\usepackage[a4paper,top=2.5cm,bottom=2.5cm,left=2.8cm,right=2.8cm,headheight=14pt]{geometry}
\usepackage{amsmath,amssymb}
\usepackage[T1]{fontenc}
\usepackage[utf8]{inputenc}
\usepackage{lmodern,microtype,fancyhdr,enumitem,hyperref,lastpage,parskip,listings}

\hypersetup{
    colorlinks=true,
    urlcolor=black,
    linkcolor=black,
    pdftitle={CSCSE31: Python Programming (Sessional Exam) -- BHU Sem III 2025-26 (NEP SEC)}
}

\pagestyle{fancy}
\fancyhf{}
\renewcommand{\headrulewidth}{0.4pt}
\renewcommand{\footrulewidth}{0.4pt}
\fancyhead[L]{\small \textit{Banaras Hindu University}}
\fancyhead[R]{\small \textit{CSCSE31: Python Programming (Sessional)}}
\fancyfoot[C]{\small Page \thepage\ of \pageref{LastPage}}

\newcommand{\pts}[1]{\hfill \textit{[#1]}}

\lstset{
  language=Python,
  basicstyle=\ttfamily\small,
  keywordstyle=\bfseries,
  showstringspaces=false,
  frame=single,
  frameround=tttt
}

\begin{document}

\begin{center}
  {\large\bfseries BANARAS HINDU UNIVERSITY}\\[2pt]
  {\normalsize B.Sc. Semester III Sessional Examination 2025-26 (NEP)}\\[6pt]
  \rule{0.8\linewidth}{0.5pt}\\[6pt]
  {\large\bfseries COMPUTER SCIENCE (SEC)}\\[4pt]
  {\large\bfseries Paper Code: CSCSE31 : Python Programming (Sessional Exam)}\\[4pt]
  {\normalsize Time: 1 Hour\hfill Maximum Marks: 20}
\end{center}

\rule{\linewidth}{0.4pt}

\smallskip

\noindent \textit{Instructions: Answer all sections as specified.}

\bigskip

\begin{center}
  \textbf{Part A --- One Word / Very Short Answer Questions}\\
  \textit{(Answer any five. Each question carries 1 mark. $5 \times 1 = 5$ marks)}
\end{center}
\begin{enumerate}[label=\textbf{\arabic*.}]
  \item Which function is used to get input from the user in Python?
  \item What is the output of \texttt{3 ** 2}?
  \item Name the keyword used to exit a loop prematurely.
  \item Write one example of a logical operator.
  \item What does the \texttt{append()} method do for a list?
  \item What data type does \texttt{\{\}} represent in Python?
  \item What is the return value of a function without a return statement?
\end{enumerate}

\bigskip

\begin{center}
  \textbf{Part B --- Short Answer Questions}\\
  \textit{(Answer any four. Each question carries 2.5 marks. $4 \times 2.5 = 10$ marks)}
\end{center}
\begin{enumerate}[label=\textbf{\arabic*.}, start=8]
  \item Explain the difference between \texttt{and}, \texttt{or}, and \texttt{not} operators with examples.
  \item Write a Python program to check whether a number entered by the user is even or odd.
  \item Write code to create a list of numbers \texttt{[1, 2, 3, 4, 5]} and print the sum of all elements.
  \item What is the difference between a list and a dictionary? Give one example of each.
  \item Define a function \texttt{greet(name)} that prints a greeting message for the given name.
\end{enumerate}

\bigskip

\begin{center}
  \textbf{Part C --- Long Answer Question}\\
  \textit{(Answer any one. Each question carries 5 marks. $1 \times 5 = 5$ marks)}
\end{center}

\noindent \textbf{Q13. (a)} Write a Python program to accept the age of a person and determine the ticket price for a movie based on age using if-elif-else (Children below 12 years: Rs. 100, Teenagers 12--17: Rs. 150, Adults 18 and above: Rs. 200). \pts{5}\\
\begin{center}\textbf{OR}\end{center}
\noindent \textbf{Q13. (b)} Write a Python function \texttt{addNum()} that may accept zero or more numbers and give sum and average of those numbers. Show an example call. \pts{5}

\bigskip
\begin{center}
  \textbf{***}
\end{center}

\end{document}
